\chapter{2015年北京卷（理）}

\section{单选题}

\begin{problem}
复数 \(\i(2 - \i) =\)（\quad）

\choices
    {\(1 + 2\i\)}
    {\(1 - 2\i\)}
    {\(-1 + 2\i\)}
    {\(-1 - 2\i\)}

\begin{answer}
A
\end{answer}

\begin{solution}
由复数的乘法及 \(\i^2=-1\)，得 \(\i(2-\i)=2\i-\i^2=1+2\i\)，所以选 A.
\end{solution}
\end{problem}

\begin{problem}
若 \(x\)，\(y\) 满足 \(\begin{cases}
x - y \leqslant 0,\\
x + y \leqslant 1,\\
x \geqslant 0,
\end{cases}\) 则 \(z = x + 2y\) 最大值为（\quad）

\choices
    {\(0\)}
    {\(1\)}
    {\( \frac{3}{2} \)}
    {\(2\)}

\begin{answer}
D
\end{answer}

\begin{solution}
由 \(x-y\leqslant0\)、\(x+y\leqslant1\) 和 \(x\geqslant0\)，可知可行域的顶点为 \((0,0)\)、\((0,1)\)、\(\left(\frac{1}{2},\frac{1}{2}\right)\). 目标函数 \(z=x+2y\) 在这些顶点处的值分别为 \(0\)、\(2\)、\(\frac{3}{2}\)，故最大值为 \(2\)，所以选 D.
\end{solution}
\end{problem}

\begin{problem}
执行如图所示的程序框图，输出的结果为（\quad）

\texfigure[max width=0.46\linewidth]{img_repaint/2015/beijing_science/q3_fig1.tex}

\choices
    {\((-2, 2)\)}
    {\((-4,0)\)}
    {\((-4, -4)\)}
    {\((0, 8)\)}

\begin{answer}
B
\end{answer}

\begin{solution}
初始时 \((x,y,k)=(1,1,0)\). 每循环一次先计算 \(s=x-y\)、\(t=x+y\)，再令 \(x=s\)、\(y=t\)，并使 \(k\) 增加 \(1\). 因而各次循环后的状态为
\((x,y,k)=(0,2,1)\)，\((-2,2,2)\)，\((-4,0,3)\).
此时满足 \(k\geqslant3\)，输出 \((-4,0)\)，所以选 B.
\end{solution}
\end{problem}

\begin{problem}
设 \(\alpha\)，\(\beta\) 是两个不同的平面，\(m\) 是直线且 \(m \subset \alpha\)，\(m \parallel \beta\) 是“\(\alpha \parallel \beta\)”的（\quad）

\choices
    {充分而不必要条件}
    {必要而不充分条件}
    {充分必要条件}
    {既不充分也不必要条件}

\begin{answer}
B
\end{answer}

\begin{solution}
若 \(\alpha\parallel\beta\)，则 \(\alpha\) 内任意直线都与平面 \(\beta\) 平行，因此 \(\alpha\parallel\beta\) 能推出 \(m\parallel\beta\)，所以前者是后者的必要条件.

反之，若 \(\alpha\) 与 \(\beta\) 相交，取 \(m\) 为 \(\alpha\) 内平行于交线 \(\alpha\cap\beta\) 的直线，则 \(m\parallel\beta\)，但 \(\alpha\) 不平行于 \(\beta\). 因此该条件不是充分条件，故选 B.
\end{solution}
\end{problem}

\begin{problem}
某三棱锥的三视图如图所示，则该三棱锥的表面积是（\quad）

\texfigure[max width=0.65\linewidth]{img_repaint/2015/beijing_science/q5_fig1.tex}

\choices
    {\(2 + \sqrt{5}\)}
    {\(4 + \sqrt{5}\)}
    {\(2 + 2\sqrt{5}\)}
    {\(5\)}

\begin{answer}
C
\end{answer}

\begin{solution}
设三棱锥顶点为 \(P\)，底面为 \(\triangle ABC\)，由三视图可还原出 \(PC\perp\) 平面 \(ABC\)，且 \(PC=1\). 设 \(D\) 是 \(AB\) 的中点，则 \(AB=2\)、\(AD=BD=1\)、\(CD=2\)，且 \(CD\perp AB\). 因而
\(S_{\triangle ABC}=\frac{1}{2}\times2\times2=2\)，\(AC=BC=\sqrt{1^2+2^2}=\sqrt5\).
又因为 \(PC\perp\) 平面 \(ABC\)，所以 \(PC\perp AC\)、\(PC\perp BC\)，从而
\[
S_{\triangle PAC}=S_{\triangle PBC}=\frac{1}{2}\times1\times\sqrt5=\frac{\sqrt5}{2}
\]
在直角三角形 \(PCD\) 中，\(PD=\sqrt{PC^2+CD^2}=\sqrt5\). 又 \(AB\perp PC\)、\(AB\perp CD\)，而 \(PC\)、\(CD\) 是平面 \(PCD\) 内相交的两条直线，所以 \(AB\perp\) 平面 \(PCD\)，从而 \(PD\perp AB\). 因此 \(S_{\triangle PAB}=\frac{1}{2}\times2\times\sqrt5=\sqrt5\). 三棱锥的表面积为 \(2+\sqrt5+\frac{\sqrt5}{2}+\frac{\sqrt5}{2}=2+2\sqrt5\)，故选 C.
\end{solution}
\end{problem}

\begin{problem}
设 \(\{a_{n}\}\) 是等差数列，下列结论中正确的是（\quad）

\choices
    {若 \( a_{1} + a_{2} > 0 \)，则 \( a_{2} + a_{3} > 0 \)}
    {若 \( a_{1} + a_{3} < 0 \)，则 \( a_{1} + a_{2} < 0 \)}
    {若 \(0 < a_{1} < a_{2}\)，则 \(a_{2} > \sqrt{a_{1}a_{3}}\)}
    {若 \(a_{1} < 0\)，则 \((a_{2} - a_{1})(a_{2} - a_{3}) > 0\)}

\begin{answer}
C
\end{answer}

\begin{solution}
设等差数列的公差为 \(d\)，则 \(a_2=a_1+d\)、\(a_3=a_1+2d\).

选项 A 不成立，例如取 \(a_1=1\)、\(d=-\frac{3}{2}\)，则 \(a_1+a_2=\frac{1}{2}>0\)，但 \(a_2+a_3=-\frac{5}{2}<0\). 选项 B 不成立，例如取 \(a_1=1\)、\(d=-\frac{6}{5}\)，则 \(a_1+a_3=-\frac{2}{5}<0\)，但 \(a_1+a_2=\frac{4}{5}>0\).

若 \(0<a_1<a_2\)，则 \(d>0\)，且
\[
a_2^2-a_1a_3=(a_1+d)^2-a_1(a_1+2d)=d^2>0
\]
由于 \(a_1a_3>0\)，故 \(a_2>\sqrt{a_1a_3}\)，选项 C 正确. 对于选项 D，有 \((a_2-a_1)(a_2-a_3)=d(-d)=-d^2\leqslant0\)，不可能大于零，故选 C.
\end{solution}
\end{problem}

\begin{problem}
如图，函数 \( f(x) \) 的图象为折线 \(ACB\)，则不等式 \( f(x) \geqslant \log_2(x + 1) \) 的解集是（\quad）

\texfigure[max width=0.5\linewidth]{img_repaint/2015/beijing_science/q7_fig1.tex}

\choices
    {\( \{x\mid-1<x\leq0\} \)}
    {\(\{x\mid -1\leqslant x\leqslant 1\}\)}
    {\(\{x\mid -1 < x\leqslant 1\}\)}
    {\(\{x\mid -1 < x\leqslant 2\}\)}

\begin{answer}
C
\end{answer}

\begin{solution}
由图可知，折线 \(ACB\) 的端点横坐标为 \(-1\)、\(2\)，且 \(C=(0,2)\). 当 \(-1<x\leqslant0\) 时，\(f(x)=2x+2=2(x+1)>0\)，而 \(0<x+1\leqslant1\) 时 \(\log_2(x+1)\leqslant0\)，所以不等式成立.

当 \(0\leqslant x\leqslant1\) 时，\(f(x)=2-x\geqslant1\)，且 \(1\leqslant x+1\leqslant2\)，故 \(\log_2(x+1)\leqslant1\)，不等式仍成立. 当 \(1<x\leqslant2\) 时，\(f(x)=2-x<1\)，而 \(\log_2(x+1)>1\)，不等式不成立. 另外，\(x=-1\) 不在对数函数定义域内，因此解集为 \(\{x\mid-1<x\leqslant1\}\)，故选 C.
\end{solution}
\end{problem}

\begin{problem}
汽车的“燃油效率”，是指汽车每消耗\(1\text{ 升}\)汽油行驶的里程，下图描述了甲，乙，丙三辆汽车在不同速度下的燃油效率情况. 下列叙述中正确的是（\quad）

\texfigure[max width=0.55\linewidth]{img_repaint/2015/beijing_science/q8_fig1.tex}

\choices
    {消耗\(1\text{ 升}\)汽油，乙车最多可行驶\(5\text{ 千米}\)}
    {以相同速度行驶相同路程，三辆车中，甲车消耗汽油最多}
    {甲车以\(80\text{ 千米/小时}\)的速度行驶\(1\text{ 小时}\)，消耗\(10\text{ 升}\)汽油}
    {某城市机动车最高限速为\(80\text{ 千米/小时}\). 相同条件下，在该市用丙车比用乙车更省油}

\begin{answer}
D
\end{answer}

\begin{solution}
燃油效率越大，在相同路程下消耗的汽油越少. 由图可知，乙车的燃油效率最大值超过 \(5\text{ 千米/升}\)，故 A 错；在相同速度和路程下，甲车的燃油效率高于丙车，故甲车不可能消耗汽油最多，B 错；甲车以 \(80\text{ 千米/小时}\) 行驶 \(1\) 小时的路程为 \(80\text{ 千米}\)，图中该速度下燃油效率为 \(10\text{ 千米/升}\)，所以消耗汽油 \(80\div10=8\text{ 升}\)，C 错.

在不超过 \(80\text{ 千米/小时}\) 的速度范围内，图中丙车的燃油效率不低于乙车，因此相同条件下丙车消耗的汽油不多于乙车，故 D 正确.
\end{solution}
\end{problem}

\section{填空题}

\begin{problem}
在 \((2 + x)^5\) 的展开式中，\(x^3\) 的系数为\fillinblank{}.

\begin{answer}
40
\end{answer}

\begin{solution}
二项式展开式的通项为 \(\mathrm C_5^k2^{5-k}x^k\). 令 \(k=3\)，得 \(x^3\) 的系数为 \(\mathrm C_5^3 2^2=10\times4=40\).
\end{solution}
\end{problem}

\begin{problem}
已知双曲线 \(\dfrac{x^2}{a^2}-y^2=1\)（\(a>0\)）的一条渐近线为 \(\sqrt{3}x+y=0\)，则 \(a=\fillinblank{}\).

\begin{answer}
\(\dfrac{\sqrt{3}}{3}\)
\end{answer}

\begin{solution}
双曲线 \(\dfrac{x^2}{a^2}-y^2=1\) 的渐近线为 \(y=\pm\dfrac{x}{a}\). 已知一条渐近线为 \(y=-\sqrt3x\)，所以 \(\dfrac{1}{a}=\sqrt3\). 由 \(a>0\)，得 \(a=\dfrac{\sqrt3}{3}\).
\end{solution}
\end{problem}

\begin{problem}
在极坐标系中，点 \(\left(2,\dfrac{\pi}{3}\right)\) 到直线 \(\rho(\cos\theta+\sqrt{3}\sin\theta)=6\) 的距离为\fillinblank{}.

\begin{answer}
1
\end{answer}

\begin{solution}
极坐标点 \(\left(2,\dfrac{\pi}{3}\right)\) 的直角坐标为 \((x,y)=(1,\sqrt3)\). 直线方程化为 \(x+\sqrt3y=6\). 因此点到直线的距离为
\[
\frac{|1+\sqrt3\cdot\sqrt3-6|}{\sqrt{1+(\sqrt3)^2}}=\frac{2}{2}=1
\]
\end{solution}
\end{problem}

\begin{problem}
在 \(\triangle ABC\) 中，\(a = 4\)，\(b = 5\)，\(c = 6\)，则 \(\frac{\sin 2A}{\sin C} = \)\fillinblank{}.

\begin{answer}
1
\end{answer}

\begin{solution}
由正弦定理，\(\dfrac{\sin A}{\sin C}=\dfrac{a}{c}=\dfrac{4}{6}=\dfrac{2}{3}\). 由余弦定理，
\[
\cos A=\frac{b^2+c^2-a^2}{2bc}=\frac{5^2+6^2-4^2}{2\times5\times6}=\frac{3}{4}
\]
所以 \(\dfrac{\sin2A}{\sin C}=2\cos A\dfrac{\sin A}{\sin C}=2\times\dfrac34\times\dfrac23=1\).
\end{solution}
\end{problem}

\begin{problem}
在 \(\triangle ABC\) 中，点 \(M\)，\(N\) 满足 \(\overrightarrow{AM} = 2\overrightarrow{MC}\)，\(\overrightarrow{BN} = \overrightarrow{NC}\). 若 \(\overrightarrow{MN} = x\overrightarrow{AB} + y\overrightarrow{AC}\)，则 \(x =\)\fillinblank{}；\(y =\)\fillinblank{}.

\begin{answer}
\(x=\dfrac{1}{2}\)，\(y=-\dfrac{1}{6}\)
\end{answer}

\begin{solution}
由 \(\overrightarrow{AM}=2\overrightarrow{MC}\)，得 \(\overrightarrow{AM}=\dfrac{2}{3}\overrightarrow{AC}\). 又由 \(\overrightarrow{BN}=\overrightarrow{NC}\)，得 \(N\) 是 \(BC\) 的中点. 以 \(A\) 为起点，用 \(\overrightarrow{AB}\)、\(\overrightarrow{AC}\) 表示位置向量，则
\(\overrightarrow{AM}=\frac23\overrightarrow{AC}\)，\(\overrightarrow{AN}=\overrightarrow{AB}+\frac12\overrightarrow{BC} =\frac12\overrightarrow{AB}+\frac12\overrightarrow{AC}\).
于是
\[
\overrightarrow{MN}=\overrightarrow{AN}-\overrightarrow{AM}
=\frac12\overrightarrow{AB}-\frac16\overrightarrow{AC}
\]
故 \(x=\dfrac12\)，\(y=-\dfrac16\).
\end{solution}
\end{problem}

\begin{problem}
函数 \(f(x)=\begin{cases}
2^x-a, & x<1,\\
4(x-a)(x-2a), & x\geqslant 1.
\end{cases}\)

\begin{enumerate}
\item 若 \(a=1\)，则 \(f(x)\) 的最小值为\fillinblank{}；
\item 若 \(f(x)\) 恰有 \(2\) 个零点，则实数 \(a\) 的取值范围是\fillinblank{}.
\end{enumerate}

\begin{answer}
\(-1\)，\(\left[\dfrac{1}{2},1\right)\cup[2,+\infty)\)
\end{answer}

\begin{solution}
\begin{enumerate}
\item 当 \(a=1\) 时，\(x<1\) 时 \(f(x)=2^x-1\)，该函数递增且下界为 \(-1\). 当 \(x\geqslant1\) 时，\(f(x)=4(x-1)(x-2)\)，其顶点为 \(\left(\frac32,-1\right)\)，且顶点在定义域内. 所以 \(f(x)\) 的最小值为 \(-1\).
\item 当 \(x<1\) 时，方程 \(f(x)=0\) 等价于 \(2^x=a\)，它有一个定义域内的根当且仅当 \(0<a<2\). 当 \(x\geqslant1\) 时，方程 \(f(x)=0\) 的根为 \(x=a\) 和 \(x=2a\)，分别满足定义域条件当且仅当 \(a\geqslant1\) 和 \(a\geqslant\frac12\).

当 \(a\leqslant0\) 时没有零点；当 \(0<a<\frac12\) 时只有第一个分支的一个零点；当 \(\frac12\leqslant a<1\) 时有第一个分支的一个零点和第二分支的一个零点；当 \(1\leqslant a<2\) 时共有三个零点. 当 \(a\geqslant2\) 时第一个分支没有零点，而第二分支的两个根均在定义域内. 因此恰有两个零点时，\(a\in\left[\frac12,1\right)\cup[2,+\infty)\).
\end{enumerate}
\end{solution}
\end{problem}

\section{解答题}

\begin{problem}
已知函数 \(f(x)=\sqrt{2}\sin\dfrac{x}{2}\cos\dfrac{x}{2}-\sqrt{2}\sin^2\dfrac{x}{2}\).

\begin{enumerate}
\item 求 \(f(x)\) 的最小正周期；
\item 求 \(f(x)\) 在区间 \(\left[-\pi,0\right]\) 上的最小值.
\end{enumerate}
\end{problem}

\begin{solution}
\begin{enumerate}
\item 利用降幂公式，得
\[
f(x)=\frac{\sqrt2}{2}\sin x-\frac{\sqrt2}{2}\left(1-\cos x\right)
=\frac{\sqrt2}{2}\left(\sin x+\cos x\right)-\frac{\sqrt2}{2}
=\sin\left(x+\frac{\pi}{4}\right)-\frac{\sqrt2}{2}
\]
因此 \(f(x)\) 的最小正周期为 \(2\pi\).
\item 当 \(x\in[-\pi,0]\) 时，\(x+\dfrac{\pi}{4}\in\left[-\dfrac{3\pi}{4},\dfrac{\pi}{4}\right]\). 该区间包含 \(-\dfrac{\pi}{2}\)，所以 \(\sin\left(x+\dfrac{\pi}{4}\right)\) 的最小值为 \(-1\)，在 \(x=-\dfrac{3\pi}{4}\) 时取得. 故 \(f(x)\) 的最小值为 \(-1-\dfrac{\sqrt2}{2}\).
\end{enumerate}
\end{solution}

\begin{problem}
\(A\)，\(B\) 两组各有 \(7\) 位病人，他们服用某种药物后的康复时间（单位：天）记录如下：

\(A\) 组：\(10\)，\(11\)，\(12\)，\(13\)，\(14\)，\(15\)，\(16\)；

\(B\) 组：\(12\)，\(13\)，\(15\)，\(16\)，\(17\)，\(14\)，\(a\).

假设所有病人的康复时间相互独立，从 \(A\)，\(B\) 两组随机各选 \(1\) 人，\(A\) 组选出的人记为甲，\(B\) 组选出的人记为乙.

\begin{enumerate}
\item 求甲的康复时间不少于 \(14\) 天的概率；
\item 如果 \(a=25\)，求甲的康复时间比乙的康复时间长的概率；
\item 当 \(a\) 为何值时，\(A\)，\(B\) 两组病人康复时间的方差相等？
\end{enumerate}
\end{problem}

\begin{solution}
\begin{enumerate}
\item A 组有 \(7\) 位病人，其中康复时间不少于 \(14\) 天的有 \(14\)、\(15\)、\(16\) 共 \(3\) 位，所以所求概率为 \(\dfrac{3}{7}\).
\item 当 \(a=25\) 时，甲、乙分别从两组中各选一人，共有 \(7\times7=49\) 种等可能结果. 按甲的康复时间列举乙的康复时间较短的结果：甲为 \(13\) 时有 \(1\) 种，甲为 \(14\) 时有 \(2\) 种，甲为 \(15\) 时有 \(3\) 种，甲为 \(16\) 时有 \(4\) 种，合计 \(10\) 种. 因此所求概率为 \(\dfrac{10}{49}\).
\item A 组的平均数为 \(13\)，方差为
\[
\frac{(10-13)^2+(11-13)^2+\cdots+(16-13)^2}{7}=\frac{28}{7}=4
\]
B 组的平均数为 \(\dfrac{a+87}{7}\)，平方和为 \(a^2+1279\)，所以其方差为
\[
\frac{a^2+1279}{7}-\left(\frac{a+87}{7}\right)^2
\]
令其等于 \(4\)，得
\[
\frac{a^2+1279}{7}-\frac{(a+87)^2}{49}=4
\]
即 \(a^2-29a+198=0\)，所以 \(a=11\) 或 \(a=18\).
\end{enumerate}
\end{solution}

\begin{problem}
如图，在四棱锥 \(A-EFCB\) 中，\(\triangle AEF\) 为等边三角形，平面 \(AEF \perp\) 平面 \(EFCB\)，\(EF \parallel BC\)，\(BC=4\)，\(EF=2a\)，\(\angle EBC=\angle FCB=60^\circ\)，\(O\) 为 \(EF\) 的中点.

\begin{enumerate}
\item 求证：\(AO \perp BE\)；
\item 求二面角 \(F-AE-B\) 的余弦值；
\item 若 \(BE \perp\) 平面 \(AOC\)，求 \(a\) 的值.
\end{enumerate}

\texfigure[max width=0.45\linewidth]{img_repaint/2015/beijing_science/q17_fig1.tex}
\end{problem}

\begin{solution}
取平面 \(EFCB\) 为 \(z=0\)，令 \(O\) 为原点，\(EF\) 在 \(x\) 轴上，取
\(E=(-a,0,0)\)，\(F=(a,0,0)\)，\(B=(-2,h,0)\)，\(C=(2,h,0)\).
由 \(\angle EBC=60^\circ\) 及 \(BC=4\)，有
\[
\cos60^\circ=\frac{2-a}{\sqrt{(2-a)^2+h^2}}
\]
从而 \(h=\sqrt3(2-a)\). 因为 \(\triangle AEF\) 为边长 \(2a\) 的等边三角形，且平面 \(AEF\perp\) 平面 \(EFCB\)，可取 \(A=(0,0,\sqrt3a)\).

\begin{enumerate}
\item 因为 \(O\) 是等边三角形 \(AEF\) 的底边中点，所以 \(AO\perp EF\). 在两平面垂直时，平面 \(AEF\) 内垂直于交线 \(EF\) 的直线 \(AO\) 垂直于平面 \(EFCB\). 又 \(BE\subset\) 平面 \(EFCB\)，所以 \(AO\perp BE\).
\item 平面 \(AEF\) 的一个外法向量可取 \(\bs n_1=(0,-1,0)\). 由
\(\overrightarrow{AE}=(-a,0,-\sqrt3a)\)，\(\overrightarrow{AB}=(-2,h,-\sqrt3a)\)
可取平面 \(AEB\) 的外法向量为
\[
\bs n_2=\overrightarrow{AE}\times\overrightarrow{AB}
=a\left(\sqrt3h,\sqrt3(2-a),-h\right)
\]
代入 \(h=\sqrt3(2-a)\)，得 \(\bs n_2=a(2-a)(3,\sqrt3,-\sqrt3)\). 因而二面角 \(F-AE-B\) 的余弦值为
\[
\frac{\bs n_1\cdot\bs n_2}{|\bs n_1||\bs n_2|}
=\frac{-\sqrt3}{\sqrt{3^2+(\sqrt3)^2+(-\sqrt3)^2}}
=-\frac{\sqrt5}{5}
\]
\item 由第（1）问知 \(BE\perp AO\). 若 \(BE\perp\) 平面 \(AOC\)，还需 \(BE\perp OC\). 由坐标得
\(\overrightarrow{BE}=(2-a,-h,0)\)，\(\overrightarrow{OC}=(2,h,0)\).
所以
\[
\overrightarrow{BE}\cdot\overrightarrow{OC}=2(2-a)-h^2=0
\]
代入 \(h^2=3(2-a)^2\)，并注意四棱锥非退化，得
\(2(2-a)-3(2-a)^2=0\)，\(qquad a=\frac43\).
\end{enumerate}
\end{solution}

\begin{problem}
已知函数 \(f(x)=\ln\dfrac{1+x}{1-x}\).

\begin{enumerate}
\item 求曲线 \(y=f(x)\) 在点 \((0,f(0))\) 处的切线方程；
\item 求证：当 \(x\in(0,1)\) 时，\(f(x)>2\left(x+\dfrac{x^3}{3}\right)\)；
\item 设实数 \(k\) 使得 \(f(x)>k\left(x+\dfrac{x^3}{3}\right)\) 对 \(x\in(0,1)\) 恒成立，求 \(k\) 的最大值.
\end{enumerate}
\end{problem}

\begin{solution}
\begin{enumerate}
\item 函数的定义域为 \((-1,1)\)，且 \(f(0)=0\). 由
\[
f'(x)=\frac{1}{1+x}+\frac{1}{1-x}=\frac{2}{1-x^2}
\]
得 \(f'(0)=2\). 因此曲线在 \((0,f(0))=(0,0)\) 处的切线方程为 \(y=2x\).
\item 令 \(g(x)=f(x)-2\left(x+\dfrac{x^3}{3}\right)\). 则 \(g(0)=0\)，且当 \(0<x<1\) 时
\[
g'(x)=\frac{2}{1-x^2}-2-2x^2
=\frac{2x^4}{1-x^2}>0
\]
所以 \(g(x)>g(0)=0\)，即 \(f(x)>2\left(x+\dfrac{x^3}{3}\right)\).
\item 记 \(h(x)=x+\dfrac{x^3}{3}\). 当 \(0<x<1\) 时，\(h(x)>0\)，第（2）问说明 \(k=2\) 符合条件. 另一方面，
\[
\lim_{x\to0^+}\frac{f(x)}{h(x)}
=\lim_{x\to0^+}\frac{f'(x)}{h'(x)}
=\frac{2}{1}=2
\]
若某个实数 \(k\) 使得对所有 \(x\in(0,1)\) 都有 \(f(x)>k h(x)\)，则因 \(h(x)>0\) 有 \(\dfrac{f(x)}{h(x)}>k\)，取极限可得 \(2\geqslant k\). 所以满足条件的 \(k\) 不超过 \(2\)，最大值为 \(2\).
\end{enumerate}
\end{solution}

\begin{problem}
已知椭圆 \(C:\dfrac{x^2}{a^2}+\dfrac{y^2}{b^2}=1\)（\(a>b>0\)）的离心率为 \(\dfrac{\sqrt{2}}{2}\)，点 \(P(0,1)\) 和点 \(A(m,n)\)（\(m\neq 0\)）都在椭圆 \(C\) 上，直线 \(PA\) 交 \(x\) 轴于点 \(M\).

\begin{enumerate}
\item 求椭圆 \(C\) 的方程，并求点 \(M\) 的坐标（用 \(m\)，\(n\) 表示）；
\item 设 \(O\) 为原点，点 \(B\) 与点 \(A\) 关于 \(x\) 轴对称，直线 \(PB\) 交 \(x\) 轴于点 \(N\). 问：\(y\) 轴上是否存在点 \(Q\)，使得 \(\angle OQM=\angle ONQ\)？若存在，求点 \(Q\) 的坐标；若不存在，说明理由.
\end{enumerate}
\end{problem}

\begin{solution}
\begin{enumerate}
\item 设椭圆的半焦距为 \(c\). 由离心率 \(e=\dfrac{c}{a}=\dfrac{\sqrt2}{2}\) 及 \(c^2=a^2-b^2\)，得 \(b^2=\dfrac{a^2}{2}\). 因为 \(P(0,1)\) 在椭圆上，所以 \(b^2=1\)，从而 \(a^2=2\). 椭圆方程为 \(\dfrac{x^2}{2}+y^2=1\).

直线 \(PA\) 可参数表示为 \((x,y)=(0,1)+t(m,n-1)\). 与 \(x\) 轴相交时 \(1+t(n-1)=0\)，故 \(t=\dfrac{1}{1-n}\)，所以
\[
M\left(\frac{m}{1-n},0\right)
\]
\item 由点 \(B\) 与点 \(A\) 关于 \(x\) 轴对称，得 \(B=(m,-n)\). 同理，直线 \(PB\) 与 \(x\) 轴的交点为
\[
N\left(\frac{m}{1+n},0\right)
\]
令 \(M=(x_M,0)\)、\(N=(x_N,0)\)、\(Q=(0,y_Q)\). 在直角三角形 \(OQM\)、\(ONQ\) 中，分别在 \(Q\)、\(N\) 处取正切，角相等的充要条件为
\[
\left|\frac{x_M}{y_Q}\right|=\left|\frac{y_Q}{x_N}\right|
\]
即 \(y_Q^2=|x_Mx_N|\). 由椭圆方程 \(\dfrac{m^2}{2}+n^2=1\)，有 \(m^2=2(1-n^2)\)，因此
\[
x_Mx_N=\frac{m^2}{1-n^2}=2
\]
所以 \(y_Q^2=2\)，确实存在两个点 \(Q\)，其坐标为 \((0,\sqrt2)\) 和 \((0,-\sqrt2)\).
\end{enumerate}
\end{solution}

\begin{problem}
已知数列 \(\{a_n\}\) 满足：\(a_1\in\N^*\)，\(a_1\leqslant 36\)，且 \(a_{n+1}=\begin{cases}
2a_n, & a_n\leqslant 18,\\
2a_n-36, & a_n>18,
\end{cases}\)（\(n=1\)，\(2\)，\(\cdots\)）. 记集合 \(M=\{a_n\mid n\in\N^*\}\).

\begin{enumerate}
\item 若 \(a_1=6\)，写出集合 \(M\) 的所有元素；
\item 若集合 \(M\) 存在一个元素是 \(3\) 的倍数，证明：\(M\) 的所有元素都是 \(3\) 的倍数；
\item 求集合 \(M\) 的元素个数的最大值.
\end{enumerate}
\end{problem}

\begin{solution}
\begin{enumerate}
\item 当 \(a_1=6\) 时，数列依次为 \(6\)，\(12\)，\(24\)，\(12\)，\(24\)，\(\ldots\)，所以 \(M=\{6,12,24\}\).
\item 对任意 \(n\)，递推式都可写成 \(a_{n+1}\equiv2a_n\pmod3\)，因为 \(36\) 是 \(3\) 的倍数. 由于 \(2\) 在模 \(3\) 意义下可逆，故
\[
3\mid a_{n+1}\Longleftrightarrow3\mid a_n
\]
若 \(M\) 中有一个元素是 \(3\) 的倍数，则该元素之前的各项和之后的各项均为 \(3\) 的倍数，所以 \(M\) 的所有元素都是 \(3\) 的倍数.
\item 将递推式看作在 \(1\)，\(2\)，\(\ldots\)，\(36\) 中进行的 \(2a_n\) 对 \(36\) 取余，其中余数 \(0\) 用 \(36\) 表示. 若 \(3\mid a_1\)，令 \(a_n=3b_n\)，则 \(b_{n+1}\equiv2b_n\pmod{12}\). 从第三项开始，\(b_3\equiv4b_1\pmod{12}\)，因而后续各项至多形成一个长度为 \(2\) 的循环；连同前两项，元素个数不超过 \(4\).

若 \(3\nmid a_1\)，则所有 \(a_n\) 都不是 \(3\) 的倍数. 从第三项开始，\(a_n\) 都是 \(4\) 的倍数，可写为 \(a_n=4c_n\)，其中 \(c_n\) 按模 \(9\) 递推为 \(c_{n+1}\equiv2c_n\pmod9\). 非三倍数的模 \(9\) 非零剩余在乘以 \(2\) 后至多经过 \(6\) 项循环，因为 \(2^6\equiv1\pmod9\). 因此从第三项开始至多有 \(6\) 个不同元素，连同前两项，\(|M|\leqslant8\).

取 \(a_1=1\)，得到
\(1\)，\(2\)，\(4\)，\(8\)，\(16\)，\(32\)，\(28\)，\(20\)，\(4\)，\(\ldots\)
其中前八项互不相同，所以 \(|M|=8\) 可以达到. 因此集合 \(M\) 的元素个数的最大值为 \(8\).
\end{enumerate}
\end{solution}
